\section{Complete-polynomial-based estimator for LBB coefficient} \label{sec:estimator}

This appendix presents the detailed derivation of the LBB coefficient estimator used in this work.
This estimator provides a quantitative measure of the stability of mixed finite element formulations just by counting the degree of freedoms (DOFs) of the approximation spaces. 
The main route of derivation follows the path proposed in Ref. \cite{wu2026}.
Unlike the consideration of volumetric locking in \cite{wu2026}, here we focus on the general LBB condition for any locking issue. 
Consider the following constraint servicing to a saddle-point problem: find $u \in U$ such that 
\begin{equation}\label{eq:constraint}
    b(p,u) = g(p), \quad  \forall p \in P
\end{equation}
where $b: U\times P \to \mathbb R$, $g: P \to \mathbb R$ are bilinear and linear forms, respectively.
For an any fixed $p\in P$, there exists a linear operator $\mathcal P: U \to P$ satisfying:
\begin{equation}\label{eq:P}
    b(p,u) = (p,\mathcal P u)
\end{equation}
where $(\bullet, \bullet): P \times P \to \mathbb R$ is the inner product in $P$.

By introducing the discrete spaces $U_h \subset U$ and $P_h \subset P$, the constraint of Eq. \eqref{eq:constraint} is approximated as: find $u_h \in U_h$ such that
\begin{equation}\label{eq:discrete_constraint}
    b(p_h,u_h) = g(p_h), \quad  \forall p_h \in P_h
\end{equation}
and $\mathcal P$'s orthogonal projection operator regarding to $P_h$, namely $\mathcal P_h: U_h \to P_h$, is defined by 
\begin{equation}\label{eq:Ph}
(p_h,\mathcal P u_h - \mathcal P_h u_h) = 0
\Rightarrow
(p_h, \mathcal P u_h) = (p_h, \mathcal P_h u_h),
\quad \forall p_h \in P_h
\end{equation}

To obtain the satisfaction solution of saddle-point problems with constraint of Eq. \eqref{eq:discrete_constraint},
the well-known LBB condition is required, i.e., there exists a positive constant $\beta > 0$, namely LBB coefficient, such that
\begin{equation}
    \inf_{p_h \in P_h} \sup_{u_h \in U_h} \frac{|b(p_h,u_h)|}{\| p \|_P\| u \|_U} \geq \beta
\end{equation}
where $\| \bullet \|_U$ and $\| \bullet \|_P$ are the norms in spaces $U$ and $P$, respectively.
To further estimate the LBB coefficient $\beta$ by the degree of freedom of $U_h$ and $P_h$, i.e. $n_u = \dim U_h$ and $n_p = \dim P_h$, we introduce the following theorem.

\begin{thm}\label{thm:estimator}
(Completed-polynomial-based estimator for LBB coefficient)
Suppose that $\mathcal P: U\rightarrow P$ is the linear operator representing the bilinear form $b$ in LBB condition and $\mathcal P_h: U_h \rightarrow P_h$ is the corresponding orthogonal projection operator regarding to $P_h$,
$U_{n_u} \subset U$ is a complete polynomial space with $\dim U_{n_u} = n_u$. The LBB coefficient $\beta$ can be estimated by:
\begin{equation}
\beta \leq \inf_{U' \subset U_{n_u}\setminus \ker \mathcal P_h \mathcal I_h} \sup_{u \in U'} \frac{\| \mathcal P u\|_P}{\| u \|_U} + O(h) 
\end{equation}
where $\mathcal I_h: U \to U_h$ is the interpolation operator with the characteristic size of the discretization $h$.
\end{thm}
\begin{pf}
Firstly, with the definition of $\mathcal P$ and $\mathcal P_h$ in Eqs. \eqref{eq:P} and \eqref{eq:Ph} respectively, the LBB condition can be rewritten as:
\begin{equation}\label{eq:proof_1}
    \begin{aligned}
        \beta &\leq \inf_{p_h \in P_h} \sup_{u_h \in U_h} \frac{|(p_h,\mathcal P u_h)|}{\| p_h \|_P\| u_h \|_U} = \inf_{p_h \in P_h} \sup_{u_h \in U_h} \frac{|(p_h,\mathcal P_h u_h)|}{\| p_h \|_P\| u_h \|_U} \\
        &\leq \inf_{p_h \in \mathrm{Im} \mathcal P_h} \sup_{u \in U_h} \frac{|(p_h,\mathcal P_h u_h)|}{\| p_h \|_P\| u_h \|_U}
    \end{aligned}
\end{equation}
where $\mathrm{Im} \mathcal P_h$ is the image space of operator $\mathcal P_h$, i.e. $\mathrm{Im} \mathcal P_h := \{p_h \in P_h |\; \exists u_h \in U_h,\; \mathcal P_h u_h = p_h \}$, $\mathrm{Im} \mathcal P_h \subset P_h$.
Since both $p_h$ and $\mathcal P_h u_h$ belong to $\mathrm{Im} \mathcal P_h$, using Cauchy-Schwarz inequality, we can further obtain that:
\begin{equation}\label{eq:proof_2}
|(\mathcal P_h u_h,p_h)| \leq \| \mathcal P_h u_h \|_P \| p_h \|_P
\end{equation}
For a given $p_h$, with introduction of a space $U' := \{ u_h \in U_h |\; \mathcal P_h u_h = p_h \}$ regarding to $p_h$, the equality in Eq. \eqref{eq:proof_1} holds when $p_h = \mathcal P_h u_h$, yields:
\begin{equation}\label{eq:proof_3}
    |(\mathcal P_h u_h,p_h)| = \| \mathcal P_h u_h \|_P \| p_h \|_P, \quad \forall u_h \in U'
\end{equation}
and, obviously, $U' \subseteq U_h\setminus \ker \mathcal P_h$, $\ker \mathcal P_h$ is the kernel space of operator $\mathcal P_h$, i.e. $\ker \mathcal P_h := \{ u_h \in U_h |\; \mathcal P_h u_h = 0 \}$.
Substituting Eq. \eqref{eq:proof_3} into Eq. \eqref{eq:proof_1} leads to:
\begin{equation}\label{eq:proof_4}
    \beta \leq \inf_{U' \subset U_h\setminus \ker \mathcal P_h} \sup_{u_h \in U'} \frac{\| \mathcal P_h u_h \|_P}{\| u_h \|_U}
\end{equation}

Secondly, the complete polynomial space $V_{n_u} \subset U$ with identical dimension to $U_h$ is introduced herein to further estimate the $\beta$, where $n_u$ is the dimension of $U_h$, i.e. $n_u = \dim U_h = \dim V_{n_u}$.
At this circumstance, for any $u \in U_{n_u}$, there exists an interpolation operator $\mathcal I_h: U \to U_h$ such that $u_h = \mathcal I_h u$.
Thus, Eq. \eqref{eq:proof_4} can be rewritten as:
\begin{equation}\label{eq:proof_5}
    \beta \leq \inf_{U' \subset U_{n_u}\setminus \ker \mathcal P_h \mathcal I_h} \sup_{u \in U'} \frac{\| \mathcal P_h \mathcal I_h u \|_P}{\| \mathcal I_h u \|_U}
\end{equation}
In accordance with the triangular inequality \cite{brenner2008} and the relationship of Eq. \eqref{eq:Ph}, we have:
\begin{equation}\label{eq:proof_6}
    \begin{aligned}
        \| \mathcal P_h \mathcal I_h u \|_P & \leq \sup_{p_h \in P_h} \frac{|(\mathcal P_h \mathcal I_h u, p_h)|}{\| p_h \|_P} = \sup_{p_h \in P_h} \frac{|(\mathcal P \mathcal I_h u, p_h)|}{\| p_h \|_P} \\
        &\leq \sup_{p_h \in P_h} \frac{| (\mathcal P u, p_h) | + | (\mathcal P u - \mathcal P \mathcal I_h u, p_h) |}{\|p_h\|_P} \\
        &= \| \mathcal P u \|_P + \|\mathcal P (u - \mathcal I_h u)\|_P
    \end{aligned}
\end{equation}
where the last term is the interpolation error of space $U_h$ with respect to operator $\mathcal P$.
Generally, particularly for finite element approximations, the interpolation error is proportional to the characteristic size of the discretization $h$ \cite{brenner2008}:
\begin{equation}\label{eq:proof_7}
    \|\mathcal P (u - \mathcal I_h u)\|_P \leq C h^k |u|_{H_k}
\end{equation}
in which $C$ is a positive constant independent of $h$, $k$ is the convergence order related to the operators $\mathcal I_h$ and $\mathcal P$. With the combination of Eqs. \eqref{eq:proof_5}, \eqref{eq:proof_6} and \eqref{eq:proof_7}, the final estimator of LBB coefficient $\beta$ in Eq. \eqref{thm:estimator} can be proved:
\begin{equation}
    \begin{aligned}
    \beta &\leq \inf_{U' \subset U_{n_u}\setminus \ker \mathcal P_h \mathcal I_h} \sup_{u \in U'} \frac{\| \mathcal P_h \mathcal I_h u \|_P}{\| \mathcal I_h u \|_U} \\
    &\leq \inf_{U' \subset U_{n_u}\setminus \ker \mathcal P_h \mathcal I_h} \sup_{u \in U'} \frac{\| \mathcal P u \|_P + \| \mathcal P (u - \mathcal I_h u) \|_P}{\| \mathcal I_h u \|_U} \\
    &\leq \inf_{U' \subset U_{n_u}\setminus \ker \mathcal P_h \mathcal I_h} \sup_{u \in U'} \frac{\| \mathcal P u\|_P}{\| u \|_U} + O(h)
    \end{aligned}
\end{equation}
\end{pf}

It should be noted that Theorem \ref{thm:estimator} inherently embeds the numbers of DOFs of approximation spaces $U_h$ and $P_h$,
while $n_u = \dim U_h$ and $n_p =\dim \mathrm{Im} \mathcal P_h = \dim (U_h \setminus \ker \mathcal P_h \mathcal I_h)$,
thus it is easy to justify the LBB condition by just counting the DOFs of approximation spaces.